\section{Task and Data}
\label{sec:data}
\subsection{The unit of prediction}
Each item $i$ has the six elements listed in Table~\ref{tab:notation}. The span identifies one occurrence even when the sentence contains several abbreviations. We keep $r_i=x_i[s_i]$, including any attached prefix, and mark the span with \texttt{[ACR]} and \texttt{[/ACR]}.

\begin{table}[t]
\centering\small
\begin{tabular}{@{}cp{0.78\columnwidth}@{}}
\toprule
Symbol & Meaning \\
\midrule
$a_i$ & Acronym type, used only for grouping, never as model input \\
$x_i$ & Sentence \\
$s_i$ & Half-open character span of the target in $x_i$ \\
$r_i$ & Raw target text, $x_i[s_i]$, with any attached prefix \\
$C_i$ & Candidate expansions (the inventory) \\
$y_i$ & Reference expansion \\
\bottomrule
\end{tabular}
\caption{Elements of an item.}
\label{tab:notation}
\end{table}

Generation predicts an expansion from $(x_i,s_i,r_i)$, without candidates. Selection picks one member of $C_i$ for the same marked sentence. Both are compared with the reference $y_i$, which is also the training label. The intended answer is a true expansion of the acronym, not just a description of what it refers to.

Formally, the two formats differ in their output space:
\begin{equation}
\hat{y}^{\mathrm{gen}}=f(x,s)\in\Sigma^{*},\qquad \hat{y}^{\mathrm{sel}}=\operatorname*{arg\,max}_{c\in C}\mathrm{score}(x,s,c),
\end{equation}
where $\Sigma^{*}$ is the set of all Hebrew strings, so generation may return any string, while selection returns one candidate in $C$.

Table~\ref{tab:example} illustrates the interface with a constructed example, not a dataset item.
\begin{table}[t]
\centering\small
\begin{tabular}{@{}p{0.22\columnwidth}p{0.71\columnwidth}@{}}
\toprule
Sentence & \heb{ח״כ נשא נאום במליאה.} \\
Target & \heb{ח״כ} (the first token) \\
Generation & A short expansion, e.g.\ \heb{חבר כנסת} \\
Selection & A letter assigned to one supplied candidate \\
\bottomrule
\end{tabular}
\caption{Illustrative input and answer formats. The sentence means ``A member of parliament gave a speech in the plenum.'' The target is shown separately here for readability; both tasks receive it marked inside the sentence. The example makes no claim about system accuracy.}
\label{tab:example}
\end{table}

\subsection{Sources and construction}
Candidate expansions come from Hebrew Wikipedia abbreviation and disambiguation pages and from Wiktionary. We merged the two sources and reviewed near-duplicate expansions by hand, which left 2,702 candidate rows for 638 acronym types. Sentences come from Hebrew Wikipedia article text and from the Knesset plenary corpus (HaifaCLGroup/KnessetCorpus). We kept only sentences of 40 to 200 characters with a clear acronym boundary, allowing attached Hebrew prefixes, and we removed sentences with glosses, explicit definitions or other acronyms. A dictionary entry does not show that a sense occurs in context, so candidate sources and sentence labels are kept separate.

Sentences are of four kinds. \emph{Natural} sentences contain the acronym as written. \emph{Substituted} sentences come from Wikipedia: a spelled-out expansion was replaced with the acronym, and the removed expression is a weak label. \emph{Deglossed} sentences are natural sentences with an adjacent explanation removed. \emph{Claude-authored} sentences were written by an LLM to add senses to types that had few. Substituted and authored sentences give labels cheaply, but they need not read like naturally abbreviated text.

\subsection{Train, development and test sets}
Table~\ref{tab:data} shows how the three sets differ in size and in the kind of sentence they contain.

\begin{table}[t]
\centering\small
\begin{tabular}{@{}lrrr@{}}
\toprule
Quantity & Train & Dev & Test \\
\midrule
Items & 2,829 & 62 & 395 \\
Acronym types & 434 & 11 & 60 \\
\midrule
Substituted Wikipedia & 2,316 & 0 & 0 \\
Natural Knesset & 319 & 62 & 267 \\
Natural Wikipedia & 50 & 0 & 41 \\
Deglossed Wikipedia & 8 & 0 & 0 \\
Claude-authored & 136 & 0 & 87 \\
\bottomrule
\end{tabular}
\caption{Items by set and by kind of sentence. Every item has at least two candidates.}
\label{tab:data}
\end{table}

\paragraph{Train.} The encoder trains on 2,829 items, which give 12,684 sentence--candidate pairs. Most items are substituted, so most labels are weak. Stored labels were reused without comprehensive reannotation.

\paragraph{Development.} Development holds 62 natural Knesset items for 11 types, and we use it only to select the encoder checkpoint. Of its labels, 42 reuse an earlier human review of the same sentence, and 20 are AI-assisted author decisions, so this is not independent annotation.

\paragraph{Test.} The test set has 395 items for 60 types, and a human reviewer went over the entire set. For the Knesset sentences, the review compared each label with the candidate list: of 1,045 mined sentences, 941 were confirmed, 100 were corrected and four were dropped. All 267 Knesset test items come from this pool (234 confirmed, 33 corrected). The 87 Claude-authored items add three sentences each to 29 types whose reviewed Knesset examples all shared one sense. A single reviewer did the work, so we have no agreement estimate.

\subsection{Separation of the sets}
\label{sec:separation}
Preparation checked item identity, target boundaries and candidate structure. These checks show structural validity, not that a label is correct. We split by acronym type. The test set shares no acronym form or sentence string with train or dev, and train and dev share none either. Documents are not fully separated: train and dev share three documents (three train and five dev items), and test shares four Knesset protocols with train (five train and 14 test items). Prefixed variants of one acronym can also fall in different sets: train has \heb{מד״ר} and dev has \heb{בד״ר}. The split therefore does not show generalization to wholly unseen acronym families.
